\section{Scope}
\label{sec:scope}%
BikeBetterPath (BBP) is a platform that supports cyclists in discovering, evaluating, and using safe and comfortable routes in urban areas. The system allows registered users to record their bike trips, manually or automatically, and to submit contributions describing obstacles, roadworks, and other conditions observed on specific street segments.

The platform defines a Contribution as a structured report that associates one or more typed Information items with a StreetSegment, possibly linked to a Trip performed by a user. A StreetSegment is the atomic infrastructure unit in BBP, representing a portion of road between two consecutive intersections, with no intermediate junctions, so that each segment can be evaluated and updated independently. A BikePath is not a static object stored in the database, but a logical route dynamically computed by the system by chaining StreetSegments that are currently considered suitable for cycling, according to aggregated user feedback and automatically detected data.

The system automates key activities such as collecting sensor data from mobile devices, inferring potential obstacles, requesting user confirmation, and periodically aggregating contributions to maintain an up-to-date representation of street conditions. This aggregation follows a merge policy that weighs Information items based on their freshness, consistency, and number of independent confirmations, so that recent and widely supported reports have more influence on the segment’s state. In particular, the system evaluates specialized BikePathSuitability Information to promote or demote StreetSegments as bike paths, highlights suitable segments on the map, and uses them as the basis for route computation. Additionally, it provides tools for users to visualize recorded trips, inspect current bike paths, and obtain personalized route recommendations that reflect their preferences and the latest crowd-sourced information.

The platform relies on external geographic data providers, such as OpenStreetMap, to support
base map visualization and user navigation across the road network. These external datasets
are used exclusively to provide the underlying cartographic representation and to assist in
spatial operations (e.g., map matching).

The internal representation of bike paths and street conditions is instead fully managed by the
system through a dedicated graph-based data model, which is incrementally built and updated
from user-generated Contributions. This internal graph acts as an overlay over the base map
and represents only those StreetSegments that have been identified, validated, and aggregated
as relevant for cycling by the system.




\section{Definitions, Acronyms, Abbreviations}
\label{sec:definition_acronyms_abbreviations}

\begin{itemize}
\item \textbf{System, Platform}: these terms are used interchangeably when referring to the system-to-be-developed (BBP).
\item \textbf{BikePath}: a route suggested by the system that is composed only of StreetSegments currently marked as suitable for cycling.
\item \textbf{StreetSegment}: an atomic portion of the road network between two intersections, which can be evaluated and updated independently. StreetSegments are identified through spatial matching against external map data and are internally represented and managed by the system.
\item \textbf{Contribution}: a user-generated or automatically generated report containing one or more Information items about a StreetSegment.
\item \textbf{Information}: an abstract fact attached to a Contribution (e.g., Obstacle, Roadwork, Warning, BikePathSuitability).
\item \textbf{BikePathSuitability}: a specialized Information used to express and aggregate the suitability of a StreetSegment as part of a bike path.
\item \textbf{Trip}: a cycling session performed by a registered user, either recorded automatically through sensors or manually specified.
\end{itemize}

\begin{table}[H]
\begin{center}
\begin{tabular}{ |l|l| }
\hline
\textbf{Acronyms} & \textbf{Definition}                              \\
\hline
RASD   & Requirements Analysis \& Specification Document   \\
\hline
DD     & Design Document   \\
\hline
API    & Application Programming Interface                  \\
\hline
REST   & REpresentational State Transfer   \\
\hline
DBMS   & DataBase Management System   \\
\hline
GPS    & Global Positioning System \\
\hline
UI     & User Interface \\
\hline
JSON   & JavaScript Object Notation  \\
\hline
HTTPS  & HyperText Transfer Protocol over Secure Socket Layer   \\
\hline
BBP    & BikeBetterPath \\
\hline
\end{tabular}
\caption{Acronyms used in the document.}
\label{tab:acronyms}%
\end{center}
\end{table}

\begin{table}[H]
\begin{center}
\begin{tabular}{ |l|l| }
\hline
\textbf{Abbreviations} & \textbf{Definition} \\
\hline
UC   & Use Case \\
\hline
SD   & Sequence Diagram  \\
\hline
\end{tabular}
\caption{Abbreviations used in the document.}
\label{tab:Abbreviations}%
\end{center}
\end{table}

\section{Revision history}
\label{sec:revision_history}%
\begin{table}[H]
\begin{center}
\begin{tabular}{ |l|l|l|}
\hline
\textbf{Revised on} & \textbf{Version}   & \textbf{Description}                           \\
\hline
19/12/2025          & 1.0   & Initial Release of the document  \\
\hline
\end{tabular}
\caption{Revision history}
\label{tab:revision_history}%
\end{center}
\end{table}

\newcounter{bib}
\setcounter{bib}{1}
\newcommand{\cbib} {\thebib\stepcounter{bib}}

\section{Reference Documents}
\label{sec:reference_documents}
\begin{itemize}
    \item{[\cbib]} Di Nitto, Rossi, Camilli, \textit{"A.Y. 2025-2026 Software Engineering 2 Requirement Engineering and Design Project"}, 2025.
    \item{[\cbib]} Matteo Camilli, Elisabetta Di Nitto, Alfonso Fuggetta, Alessandro Margara, Matteo G. Rossi, Damian A. Tamburri, \textit{Software Engineering: a Structured Journey from Inception to Delivery}, Bocconi University Press, 2025, ISBN: 9791280623508

    \item{[\cbib]} \href{https://learn.microsoft.com/en-us/previous-versions/msp-n-p/dn589800(v=pandp.10)#eventual-consistency}{\textcolor{blue}{\textit{Microsoft Learn, Data Consistency primer}}}
    \label{document: data_consistency}
    
    \item{[\cbib]} \href{https://microservices.io/patterns/apigateway.html}{\textcolor{blue}{\textit{Chris Richardson, Microservices: API Gateway Pattern}}}
    \label{document: api_gateway}
    

    
    \item{[\cbib]} \href{https://learn.microsoft.com/en-us/previous-versions/msp-n-p/dn589795(v=pandp.10)#functional-partitioning}{\textcolor{blue}{\textit{Microsoft Learn, Data Partitioning guidance}}}
    \label{document: data_partitioning}
    
    \item{[\cbib]} \href{https://learn.microsoft.com/en-us/dotnet/architecture/cloud-native/relational-vs-nosql-data#the-cap-theorem}{\textcolor{blue}{\textit{Microsoft Learn, Relational vs. NoSQL data}}}
    \label{document: relational_nosql_data}
    


    \item{[\cbib]} \href{https://wiki.openstreetmap.org/}{\textcolor{blue}{\textit{OSM documentation}}}
    \label{document: OSM}

    \item{[\cbib]} \href{https://wiki.openstreetmap.org/wiki/Overpass_API}{\textcolor{blue}{\textit{Overpass API}}}
    \label{document: Overpass_API}

    \item{[\cbib]} \href{https://learn.microsoft.com/en-us/dotnet/architecture/cloud-native/distributed-data#database-per-microservice-why}{\textcolor{blue}{\textit{Microsoft Learn - Cloud-native data patterns}}}
    \label{document: cloud_native_data_patters}

\end{itemize}
\end{itemize}

\section{Overview}
\subsection{Document Structure}
\label{sec:doc_structure}%
This document is composed of seven sections:
\begin{itemize}
    \item{} \nth{1} Chapter - Introduction
    \item{} \nth{2} Chapter - Architectural Design
    \item{} \nth{3} Chapter - User Interface Design
    \item{} \nth{4} Chapter - Requirements Traceability
    \item{} \nth{5} Chapter - Implementation, Integration and Test Plan
    \item{} \nth{6} Chapter - Effort Spent
    \item{} \nth{7} Chapter - References
\end{itemize}
